%% ma: L10-B13-0116
%% lop: 10
%% chuong: 5
%% bai: 13
%% dang: 10.13.11
%% loai: TLN
%% muc: VD
%% dap_an: "182"
%% nguon: "chuyen-de-cac-so-dac-trung-cua-mau-so-lieu-khong-ghep-nhom-toan-10-chuong-trinh-moi.pdf (GV. Phan Nhật Linh) (P03), Bài 02 Dạng 2 Phần III Câu 7"
%% kien_thuc: "Lớp 10 Bài 13 (trung vị, tứ phân vị)"
%% trang_thai: chinh-thuc
%% ghi_chu: "Tự duyệt theo yêu cầu giáo viên 22/10/2026."
\begin{ex}
Điểm điều tra về chất lượng một sản phẩm mới (thang điểm $100$) như sau:
\begin{center}
\begin{tabular}{|c|c|c|c|c|c|c|c|c|c|c|c|c|}\hline
$80$ & $65$ & $51$ & $48$ & $45$ & $61$ & $30$ & $35$ & $84$ & $83$ & $60$ & $58$ & $75$\\\hline
$72$ & $68$ & $39$ & $41$ & $54$ & $61$ & $72$ & $75$ & $72$ & $61$ & $50$ & $65$ & \\\hline
\end{tabular}
\end{center}
Tính tổng của trung vị, tứ phân vị dưới và tứ phân vị trên của mẫu số liệu.
\shortans{182}
\loigiai{Sắp xếp $25$ giá trị: $30$; $35$; $39$; $41$; $45$; $48$; $50$; $51$; $54$; $58$; $60$; $61$; $61$; $61$; $65$; $65$; $68$; $72$; $72$; $72$; $75$; $75$; $80$; $83$; $84$.

$M_e=x_{13}=61$; nửa dưới gồm $12$ giá trị đầu nên $Q_1=\dfrac{x_6+x_7}2=\dfrac{48+50}2=49$; nửa trên gồm $12$ giá trị cuối nên $Q_3=\dfrac{x_{19}+x_{20}}2=72$.

Tổng: $61+49+72=182$.}
\end{ex}
